\chapter{Problemática}
\section{Problema}

La transformación digital en las organizaciones enfrenta múltiples barreras que dificultan su implementación efectiva. En primer lugar, la insuficiente atención a la cultura organizacional y la ausencia de una conceptualización clara de la cultura digital limitan el éxito de los procesos de digitalización, ya que las empresas no desarrollan las capacidades necesarias para adaptarse a los cambios tecnológicos \parencite{grover2022}. Asimismo, las pequeñas y medianas empresas (PYMES) presentan dificultades relacionadas con la escasez de recursos económicos y el desconocimiento de los factores críticos para una transformación digital exitosa, lo que restringe su capacidad de adaptación a entornos cada vez más digitalizados \parencite{jewapatarakul2024}.

Además, la adopción de tecnologías digitales se ve afectada por limitaciones financieras, bajos niveles de productividad y una brecha tecnológica respecto a las grandes empresas, factores que reducen las posibilidades de inversión en innovación y modernización de procesos \parencite{clemente2024}. De igual manera, existe una limitada evidencia práctica sobre los elementos que favorecen la transformación digital en las PYMES, lo que dificulta la formulación de estrategias efectivas para su implementación \parencite{hoang2024}.

Por otra parte, la incorporación de tecnologías de automatización presenta desafíos relacionados con la disponibilidad de recursos, la capacitación del personal y la protección de la información. Mientras las PYMES enfrentan restricciones económicas para adoptar innovaciones tecnológicas, las grandes organizaciones deben afrontar la complejidad de sus estructuras organizacionales y la necesidad de fortalecer la seguridad de sus datos \parencite{hurtado2024}.


Las organizaciones actualmente enfrentan dificultades para optimizar y gestionar eficientemente sus procesos debido a la implementación inadecuada de soluciones tecnológicas y a la falta de estrategias integrales de transformación digital. Esta situación provoca que muchas empresas adapten sus procedimientos operativos a las limitaciones de los sistemas informáticos disponibles, en lugar de seleccionar tecnologías alineadas con sus necesidades reales y objetivos organizacionales \parencite{munoz2023}.

Asimismo, la creciente complejidad de los entornos empresariales exige mecanismos de mejora continua que permitan incrementar la eficiencia operativa y responder de manera efectiva a las demandas del mercado. Sin embargo, la ausencia de herramientas de automatización y de metodologías adecuadas para la gestión de procesos limita la capacidad de las organizaciones para optimizar sus actividades y mejorar su productividad \parencite{polyvyanyy2021}.


Las limitaciones en los procesos de transformación digital y automatización generan diversas consecuencias para las organizaciones. Entre las principales se encuentran la disminución de la eficiencia operativa, el incremento de costos, la duplicidad de actividades y la reducción de la competitividad empresarial. La falta de automatización y optimización de procesos dificulta además la toma de decisiones basada en información confiable y oportuna, afectando el desempeño general de las empresas \parencite{quispe2025}.

Adicionalmente, la implementación de tecnologías que no responden a las necesidades reales de la organización provoca ineficiencias operativas, desaprovechamiento de recursos y dificultades para alcanzar mayores niveles de productividad \parencite{munoz2023}. Como resultado, las empresas enfrentan mayores obstáculos para adaptarse a los cambios tecnológicos y mantener ventajas competitivas en mercados cada vez más dinámicos.



La implementación exitosa de la transformación digital requiere no solo la incorporación de tecnologías, sino también el fortalecimiento de una cultura organizacional orientada al cambio. En este sentido, Grover et al. proponen un marco conceptual que define la cultura digital como el conjunto de recursos y prácticas que permiten a las organizaciones transformar la creación de valor mediante el uso de tecnologías \parencite{grover2022}. De manera complementaria,  el éxito de la transformación digital en las PYMES depende del fortalecimiento de la cultura digital, la adquisición de conocimiento externo y la preparación organizacional \parencite{jewapatarakul2024}.

Asimismo, la superación de las limitaciones que enfrentan las PYMES en los procesos de digitalización requiere estrategias orientadas al desarrollo de capacidades organizacionales. La formación académica de los directivos, la internacionalización y el crecimiento empresarial facilitan la adopción tecnológica \parencite{clemente2024}. De igual manera, se propone un enfoque integral basado en el marco TOE (Tecnología, Organización y Entorno), enfatizando la relevancia de la ciberseguridad, el respaldo de la alta dirección y las alianzas estratégicas para impulsar la transformación digital \parencite{hoang2024}.



Considerando las problemáticas identificadas en los procesos de transformación digital y gestión empresarial, la presente investigación toma como caso de estudio a la fábrica Confort Muebles. Este proyecto propone el desarrollo de una aplicación web que permita gestionar las proformas junto con sus respectivos detalles, controlar las liquidaciones realizadas, administrar el inventario de productos y materiales, y gestionar la información de los clientes. La implementación de esta solución tecnológica busca optimizar los procesos administrativos y comerciales de la empresa mediante la automatización de tareas, la centralización de la información y el fortalecimiento de los mecanismos de control y seguimiento.

En la fábrica Confort Muebles se ha identificado una problemática relacionada con la deficiente gestión de proformas, debido principalmente al uso de registros manuales para el manejo de la información. Esta situación genera múltiples inconvenientes, como errores en los datos, duplicidad de documentos, pérdida de información y dificultad para realizar un seguimiento adecuado a los clientes.

Asimismo, la falta de un formato estandarizado para las cotizaciones y liquidaciones provoca inconsistencias en la información, afectando la imagen de la fábrica y la confianza de los clientes. A esto se suma la ausencia de una base de datos centralizada, lo que dificulta el acceso rápido a la información y retrasa los procesos administrativos y comerciales.

Otro aspecto importante es la limitada supervisión de los documentos generados, lo cual impide un adecuado control por parte de los administradores. Como consecuencia, se presentan pérdidas de ventas, retrasos en la atención al cliente y desperdicio de tiempo en la búsqueda de información.

La problemática identificada evidencia la necesidad de fortalecer los procesos administrativos mediante la incorporación de herramientas tecnológicas que permitan gestionar de manera eficiente la información generada durante las actividades comerciales. Actualmente, la dependencia de registros manuales incrementa la probabilidad de errores humanos y limita la capacidad de la organización para mantener un control adecuado sobre las operaciones relacionadas con las proformas, liquidaciones, inventario y clientes.

Además, la inexistencia de un sistema integrado de información dificulta la trazabilidad de los procesos comerciales, ya que los datos se encuentran dispersos en documentos físicos o archivos independientes, impidiendo disponer de información actualizada y confiable para la toma de decisiones. Esta situación afecta directamente la capacidad de respuesta de la empresa frente a los requerimientos de los clientes, generando retrasos en la elaboración de cotizaciones, inconsistencias en los registros y dificultades para realizar un seguimiento oportuno de las ventas efectuadas.


Asimismo, la administración de la información de los clientes presenta restricciones debido a la falta de una base de datos centralizada que permita almacenar, consultar y actualizar los registros de manera eficiente. Como consecuencia, se dificulta la identificación del historial de compras, el seguimiento de proformas emitidas y la gestión de oportunidades comerciales, limitando la posibilidad de fortalecer las relaciones con los clientes y mejorar la calidad del servicio ofrecido.

En este contexto, la transformación digital se presenta como una alternativa para optimizar los procesos operativos y administrativos de la fábrica. La incorporación de soluciones tecnológicas permite disponer de información en tiempo real, mejorar la coordinación de las actividades internas, reducir errores en el manejo de datos y fortalecer la capacidad de gestión empresarial mediante la automatización de actividades rutinarias.

Por ello, la presente investigación propone el desarrollo de una aplicación web orientada a la gestión integral de proformas, liquidaciones, inventario y clientes en la fábrica Confort Muebles. La solución tecnológica permitirá centralizar la información en una base de datos única, automatizar los procesos de registro y consulta, estandarizar la documentación generada y proporcionar mecanismos de control que faciliten la supervisión de las actividades comerciales y administrativas.

La implementación de esta aplicación web busca contribuir a la reducción de errores asociados a los procesos manuales, mejorar la organización de la información, agilizar la atención al cliente y optimizar el control de inventario. De igual manera, permitirá disponer de reportes y consultas oportunas que apoyen la toma de decisiones estratégicas, favoreciendo la eficiencia operativa y la competitividad de la empresa en un entorno empresarial cada vez más digitalizado.

En consecuencia, surge la necesidad de analizar los procesos actuales de gestión de información en la fábrica Confort Muebles e identificar los requerimientos funcionales y tecnológicos necesarios para el diseño e implementación de una solución informática que contribuya a resolver las deficiencias existentes. De esta manera, la investigación busca aportar una herramienta tecnológica que apoye la modernización de la empresa y fortalezca sus capacidades para enfrentar los desafíos de la transformación digital.
